\documentclass[11pt,a4paper]{article}
\usepackage[utf8]{inputenc}
\usepackage[english]{babel}
\usepackage{geometry}
\geometry{a4paper, margin=1in}
\usepackage{amsmath}
\usepackage{listings}
\usepackage{xcolor}
\usepackage{hyperref}
\usepackage{fancyhdr}
\usepackage{titlesec}
\usepackage{enumitem}

% Python code styling
\lstdefinestyle{pythonstyle}{
    language=Python,
    basicstyle=\ttfamily\small,
    keywordstyle=\color{blue}\bfseries,
    commentstyle=\color{green!60!black}\itshape,
    stringstyle=\color{red},
    showstringspaces=false,
    breaklines=true,
    frame=single,
    numbers=left,
    numberstyle=\tiny\color{gray},
    backgroundcolor=\color{gray!10},
    captionpos=b
}

\lstset{style=pythonstyle}

% Header and footer
\pagestyle{fancy}
\fancyhf{}
\rhead{Plotter Class Documentation}
\lhead{Crystallographic Plotting Library}
\rfoot{Page \thepage}

% Title formatting
\titleformat{\section}{\Large\bfseries}{\thesection}{1em}{}
\titleformat{\subsection}{\large\bfseries}{\thesubsection}{1em}{}
\titleformat{\subsubsection}{\normalsize\bfseries}{\thesubsubsection}{1em}{}

\title{\textbf{Plotter Class Documentation}\\
\large Crystallographic Plotting Library Module}
\author{Detailed Class Documentation}
\date{\today}

\begin{document}

\maketitle
\tableofcontents
\newpage

\section{Overview}

The \texttt{plotter} class provides a comprehensive framework for creating publication-quality crystallographic plots including stereographic projections, pole figures, and orientation density maps.

\subsection{Key Features}
\begin{itemize}[leftmargin=*]
    \item Stereographic projections (Schmidt/Wulff nets)
    \item Pole figures with density contours
    \item Scatter plots with crystallographic indexing
    \item Interactive data annotations
    \item Customizable colormaps and color scales
    \item Support for multiple crystal systems and symmetry operations
\end{itemize}

\subsection{Attribute Categories}
\begin{itemize}[leftmargin=*]
    \item \textbf{Colorbar}: Color scale display and formatting
    \item \textbf{Figure}: Overall figure layout and titles
    \item \textbf{Colormap}: Data-to-color mapping and contours
    \item \textbf{Projection}: Stereographic projection type and settings
    \item \textbf{Axis}: Plot styling, ticks, labels
    \item \textbf{Crystal}: Crystal structure and symmetry operations
    \item \textbf{Save}: File output settings
    \item \textbf{Scatter}: Scatter plot data and styling
    \item \textbf{Interactive}: Mouse-over annotations
\end{itemize}

\section{Class Initialization}

\subsection{Method: \texttt{\_\_init\_\_()}}

\subsubsection{Description}
Initialize the plotter with default attributes for all plotting categories. All attributes can be modified using the \texttt{setAttributes()} method.

\subsubsection{Input}
\begin{itemize}[leftmargin=*]
    \item None
\end{itemize}

\subsubsection{Output}
\begin{itemize}[leftmargin=*]
    \item None (initializes plotter instance)
\end{itemize}

\subsubsection{Usage Example}
\begin{lstlisting}[language=Python, caption=Basic initialization]
import numpy as np
from plotlib import plotter

# Create plotter instance
p = plotter()

# Access default values
print("Default colormap:", p.cmap)
print("Default projection sphere:", p.sphere)

# Modify attributes
p.setAttributes(cmap='viridis', sphere='half')
\end{lstlisting}

\section{Core Methods}

\subsection{Method: \texttt{setAttributes(**kwargs)}}

\subsubsection{Description}
Set or update plotter attributes dynamically. This method allows flexible configuration of any plotter attribute without needing to know the internal structure.

\subsubsection{Input}
\begin{itemize}[leftmargin=*]
    \item \texttt{**kwargs}: keyword arguments -- Any attribute name and value pairs
\end{itemize}

\subsubsection{Output}
\begin{itemize}[leftmargin=*]
    \item None (updates instance attributes)
\end{itemize}

\subsubsection{Usage Example}
\begin{lstlisting}[language=Python, caption=Setting attributes]
p = plotter()

# Set single attribute
p.setAttributes(cmap='viridis')

# Set multiple attributes
p.setAttributes(
    sphere='half',
    ProjType='equalarea',
    figsize=(8, 8),
    dirs=[[1,0,0], [1,1,0]],
    cmap='jet'
)

# Verify changes
print("Sphere:", p.sphere)
print("Colormap:", p.cmap)
\end{lstlisting}

\subsection{Method: \texttt{plotProj(**kwargs)}}

\subsubsection{Description}
Create a stereographic projection (pole figure) with appropriate grid. Sets up the figure, axes, and projection type (equal-area Schmidt net or equal-angle Wulff net). Handles full sphere, hemisphere, or standard stereographic triangle projections.

\subsubsection{Input}
\begin{itemize}[leftmargin=*]
    \item \texttt{ProjType}: str -- \texttt{'equalarea'} (Schmidt) or \texttt{'equalangle'} (Wulff)
    \item \texttt{sphere}: str -- \texttt{'full'}, \texttt{'half'}, or \texttt{'triangle'}
    \item \texttt{figsize}: tuple -- Figure size (width, height) in inches
    \item \texttt{stereogrid}: bool -- Display stereographic grid
    \item \texttt{stereoresolution}: int -- Grid resolution
    \item \texttt{stereomesh}: bool -- Display mesh
    \item \texttt{ax}: matplotlib axis -- Existing axis to plot on (optional)
    \item \texttt{fig}: matplotlib figure -- Existing figure (optional)
\end{itemize}

\subsubsection{Output}
\begin{itemize}[leftmargin=*]
    \item None (creates/modifies \texttt{self.fig} and \texttt{self.ax})
\end{itemize}

\subsubsection{Usage Examples}

\paragraph{Example 1: Equal-area projection, upper hemisphere}
\begin{lstlisting}[language=Python, caption=Equal-area hemisphere projection]
import matplotlib.pyplot as plt

# Equal-area projection, upper hemisphere
p = plotter()
p.plotProj(ProjType='equalarea', sphere='half', figsize=(6, 6))
plt.show()
\end{lstlisting}

\paragraph{Example 2: Equal-angle projection, full sphere}
\begin{lstlisting}[language=Python, caption=Equal-angle full sphere projection]
# Equal-angle projection, full sphere
p2 = plotter()
p2.plotProj(ProjType='equalangle', sphere='full')
plt.show()
\end{lstlisting}

\paragraph{Example 3: Standard stereographic triangle}
\begin{lstlisting}[language=Python, caption=Stereographic triangle for cubic systems]
# Standard stereographic triangle (for cubic)
p3 = plotter()
p3.plotProj(
    ProjType='equalarea',
    sphere='triangle',
    stereogrid=True,
    stereoresolution=10
)
plt.show()
\end{lstlisting}

\subsection{Method: \texttt{getScales(vmcbar, numticks=None, ticks=None, tickslabels=None, geq=False, leq=False, cmapbins=100, cmapbinsmult=None)}}

\subsubsection{Description}
Generate color scale parameters including ticks, labels, and colormap. Creates a gradient colormap from white $\rightarrow$ blue $\rightarrow$ green $\rightarrow$ yellow $\rightarrow$ dark red with customizable tick positions and labels.

\subsubsection{Input}
\begin{itemize}[leftmargin=*]
    \item \texttt{vmcbar}: list/array [min, max] -- Value range for colorbar
    \item \texttt{numticks}: int -- Number of tick marks (optional, auto-computed if None)
    \item \texttt{ticks}: array -- Explicit tick positions (optional)
    \item \texttt{tickslabels}: list of str -- Custom tick labels (optional)
    \item \texttt{geq}: bool -- Add $\geq$ symbol to maximum tick (default: False)
    \item \texttt{leq}: bool -- Add $\leq$ symbol to minimum tick (default: False)
    \item \texttt{cmapbins}: int -- Number of colormap bins (default: 100)
    \item \texttt{cmapbinsmult}: int -- Multiplier for bins based on ticks (optional)
\end{itemize}

\subsubsection{Output}
Dictionary containing:
\begin{itemize}[leftmargin=*]
    \item \texttt{'tickslabels'}: list of str -- Formatted tick labels
    \item \texttt{'vm'}: list [min, max] -- Colormap value range
    \item \texttt{'vmbar'}: list [min, max] -- Colorbar value range
    \item \texttt{'cmap'}: matplotlib colormap -- Generated colormap
    \item \texttt{'ticks'}: array -- Tick positions
\end{itemize}

\subsubsection{Usage Examples}

\paragraph{Example 1: Basic scale}
\begin{lstlisting}[language=Python, caption=Basic color scale]
p = plotter()

# Basic scale from 0 to 10
scales = p.getScales(vmcbar=[0, 10])
print("Tick labels:", scales['tickslabels'])
print("Colormap range:", scales['vm'])
\end{lstlisting}

\paragraph{Example 2: Custom ticks with inequality symbols}
\begin{lstlisting}[language=Python, caption=Custom ticks with inequality symbols]
import numpy as np

scales2 = p.getScales(
    vmcbar=[0, 100],
    ticks=np.array([0, 25, 50, 75, 100]),
    geq=True,  # Add >= to maximum
    leq=True   # Add <= to minimum
)
print("Labels:", scales2['tickslabels'])
# Output: ['<=0', '25', '50', '75', '>=100']
\end{lstlisting}

\paragraph{Example 3: High-resolution colormap}
\begin{lstlisting}[language=Python, caption=High-resolution smooth colormap]
scales3 = p.getScales(
    vmcbar=[0, 1],
    cmapbins=1000  # Smooth gradient
)
\end{lstlisting}

\subsection{Method: \texttt{getFigparam(fontsize=None, save=False, phase='A,', figsize=None, **kwargs)}}

\subsubsection{Description}
Get figure parameter dictionary for saving high-quality figures.

\subsubsection{Input}
\begin{itemize}[leftmargin=*]
    \item \texttt{fontsize}: float -- Font size for text (optional)
    \item \texttt{save}: bool -- If True, use smaller figure size (default: False)
    \item \texttt{phase}: str -- Phase identifier (\texttt{'A'} or other) affects default size
    \item \texttt{figsize}: tuple -- Custom figure size (width, height) in inches
    \item \texttt{**kwargs}: Additional parameters to override defaults
\end{itemize}

\subsubsection{Output}
Dictionary containing figure parameters for \texttt{plt.savefig()}:
\begin{itemize}[leftmargin=*]
    \item \texttt{'dpi'}: int -- Resolution (default: 300)
    \item \texttt{'facecolor'}: str -- Background color
    \item \texttt{'edgecolor'}: str -- Edge color
    \item \texttt{'orientation'}: str -- \texttt{'portrait'} or \texttt{'landscape'}
    \item \texttt{'format'}: str -- Image format
    \item \texttt{'transparent'}: bool -- Transparent background
    \item \texttt{'bbox\_inches'}: str -- Bounding box setting
    \item \texttt{'pad\_inches'}: float -- Padding around figure
\end{itemize}

\subsubsection{Usage Examples}

\paragraph{Example 1: Default parameters}
\begin{lstlisting}[language=Python, caption=Get default figure parameters]
p = plotter()

# Get default parameters
params = p.getFigparam()
print("DPI:", params['dpi'])
\end{lstlisting}

\paragraph{Example 2: High-resolution PDF}
\begin{lstlisting}[language=Python, caption=High-resolution PDF parameters]
params_pdf = p.getFigparam(
    save=True,
    format='pdf',
    dpi=600
)
\end{lstlisting}

\paragraph{Example 3: Use with savefig}
\begin{lstlisting}[language=Python, caption=Using parameters with matplotlib]
import matplotlib.pyplot as plt

fig, ax = plt.subplots()
ax.plot([1, 2, 3], [1, 4, 9])
fig.savefig('output.png', **params)
\end{lstlisting}

\subsection{Method: \texttt{figsave(**kwargs)}}

\subsubsection{Description}
Save figure to file(s) with optional cropping. Supports multiple image formats and automatic cropping of whitespace. Can save to multiple formats simultaneously.

\subsubsection{Input}
\begin{itemize}[leftmargin=*]
    \item \texttt{fname}: str or list -- Filename(s) to save
    \item \texttt{imformats}: list of str -- Image formats [\texttt{'png'}, \texttt{'pdf'}, \texttt{'svg'}, etc.]
    \item \texttt{crop}: bool -- Auto-crop whitespace (requires wand library)
    \item \texttt{figparam}: dict -- Figure parameters (see \texttt{getFigparam})
\end{itemize}

\subsubsection{Output}
\begin{itemize}[leftmargin=*]
    \item None (saves file(s) to disk)
\end{itemize}

\subsubsection{Usage Examples}

\paragraph{Example 1: Save single format}
\begin{lstlisting}[language=Python, caption=Save to single format]
p = plotter()
p.plotProj(ProjType='equalarea')

# Save single format
p.figsave(fname='pole_figure.png')
\end{lstlisting}

\paragraph{Example 2: Save multiple formats}
\begin{lstlisting}[language=Python, caption=Save to multiple formats simultaneously]
# Save multiple formats
p.figsave(
    fname='pole_figure.png',
    imformats=['png', 'pdf', 'svg']
)
# Creates: pole_figure.png, pole_figure.pdf, pole_figure.svg
\end{lstlisting}

\paragraph{Example 3: Save with cropping}
\begin{lstlisting}[language=Python, caption=Save with automatic whitespace cropping]
# Save with cropping
p.figsave(
    fname='cropped_figure.png',
    crop=True
)
\end{lstlisting}

\paragraph{Example 4: High DPI save}
\begin{lstlisting}[language=Python, caption=High resolution output]
# High DPI save
p.figparam['dpi'] = 600
p.figsave(fname='high_res.png')
\end{lstlisting}

\subsection{Method: \texttt{figsaveproc(fname, **kwargs)}}

\subsubsection{Description}
Process and save figure to file. Internal method called by \texttt{figsave()} to handle actual file writing.

\subsubsection{Input}
\begin{itemize}[leftmargin=*]
    \item \texttt{fname}: str -- Filename to save
    \item \texttt{**kwargs}: Additional parameters to update
\end{itemize}

\subsubsection{Output}
\begin{itemize}[leftmargin=*]
    \item None (saves file to disk)
\end{itemize}

\subsubsection{Usage Example}
\begin{lstlisting}[language=Python, caption=Direct file saving (typically called internally)]
# Typically called internally by figsave()
# But can be used directly:
p = plotter()
p.plotProj()
p.figsaveproc('output.png')
\end{lstlisting}

\section{Complete Usage Examples}

\subsection{Basic Pole Figure}

\begin{lstlisting}[language=Python, caption=Creating a basic pole figure]
import numpy as np
from plotlib import plotter

# Create plotter instance
p = plotter()

# Setup projection (equal-area, upper hemisphere)
p.plotProj(ProjType='equalarea', sphere='half')

# Plot crystal directions
p.setAttributes(
    dirs=[[1,0,0], [0,1,0], [0,0,1]],
    norms=[[1,1,1]]
)
p.plotDirsNorms()

# Save figure
p.figsave(fname='pole_figure.png')
\end{lstlisting}

\subsection{Density Contour Plot}

\begin{lstlisting}[language=Python, caption=Creating an orientation density plot]
import numpy as np
from plotlib import plotter

# Random orientations
orientations = np.random.randn(1000, 3, 3)

# Create density plot
p2 = plotter()
p2.setAttributes(oris=orientations, ProjType='equalarea')
p2.plotProj()
p2.plotColormap()
p2.plotColorbar()

# Save with multiple formats
p2.figsave(fname='density_plot.png', imformats=['png', 'pdf'])
\end{lstlisting}

\subsection{Advanced Multi-Panel Figure}

\begin{lstlisting}[language=Python, caption=Creating a multi-panel crystallographic figure]
import matplotlib.pyplot as plt
import numpy as np
from plotlib import plotter

# Create 2x2 subplot
fig, axes = plt.subplots(2, 2, figsize=(12, 12))

# Create multiple pole figures
for i, ax in enumerate(axes.flat):
    p = plotter()
    p.plotProj(
        ProjType='equalarea',
        sphere='half',
        ax=ax,
        fig=fig
    )
    
    # Different directions for each subplot
    dirs = [[i+1, 0, 0], [0, i+1, 0], [0, 0, i+1]]
    p.setAttributes(dirs=dirs)
    p.plotDirsNorms()

plt.tight_layout()
plt.savefig('multi_panel_figure.png', dpi=300)
\end{lstlisting}

\section{Standalone Utility Functions}

\subsection{Function: \texttt{get\_cmap(colors, nbins=1000, name='my\_cmap')}}

\subsubsection{Description}
Create a custom colormap from a list of colors with smooth gradients.

\subsubsection{Input}
\begin{itemize}[leftmargin=*]
    \item \texttt{colors}: list of tuples/colors -- Colors to interpolate between (RGB tuples or named colors)
    \item \texttt{nbins}: int -- Number of discrete bins in colormap (default: 1000)
    \item \texttt{name}: str -- Name for the colormap (default: \texttt{'my\_cmap'})
\end{itemize}

\subsubsection{Output}
\begin{itemize}[leftmargin=*]
    \item \texttt{cmap}: matplotlib.colors.LinearSegmentedColormap -- Custom colormap
\end{itemize}

\subsubsection{Usage Example}
\begin{lstlisting}[language=Python, caption=Creating custom colormaps]
import matplotlib.pyplot as plt
import numpy as np
from plotlib import get_cmap

# Create colormap from white to red to black
colors = [(1, 1, 1), (1, 0, 0), (0, 0, 0)]
cmap = get_cmap(colors, nbins=256)

# Use in plot
data = np.random.rand(10, 10)
plt.imshow(data, cmap=cmap)
plt.colorbar()
plt.show()

# Multi-color gradient
colors_multi = [
    (0, 0, 1),      # Blue
    (0, 1, 1),      # Cyan
    (0, 1, 0),      # Green
    (1, 1, 0),      # Yellow
    (1, 0, 0)       # Red
]
cmap_rainbow = get_cmap(colors_multi, nbins=500)
\end{lstlisting}

\subsection{Function: \texttt{get\_colors(values, cmap, vmin=None, vmax=None, cmin=0, cmax=None, to255=True, nancolor=[0, 0, 0, 1])}}

\subsubsection{Description}
Map data values to colors using a colormap.

\subsubsection{Input}
\begin{itemize}[leftmargin=*]
    \item \texttt{values}: numpy array -- Data values to map to colors
    \item \texttt{cmap}: matplotlib colormap -- Colormap to use
    \item \texttt{vmin}: float -- Minimum data value (default: min of values)
    \item \texttt{vmax}: float -- Maximum data value (default: max of values)
    \item \texttt{cmin}: int -- Minimum colormap index (default: 0)
    \item \texttt{cmax}: int -- Maximum colormap index (default: \texttt{cmap.N})
    \item \texttt{to255}: bool -- Scale colors to 0-255 range (default: True)
    \item \texttt{nancolor}: list -- RGBA color for NaN values (default: black)
\end{itemize}

\subsubsection{Output}
\begin{itemize}[leftmargin=*]
    \item \texttt{Colors}: numpy array -- Array of colors (RGBA), shape: (\texttt{values.shape}, 4)
\end{itemize}

\subsubsection{Usage Example}
\begin{lstlisting}[language=Python, caption=Mapping values to colors]
import numpy as np
import matplotlib.pyplot as plt
from plotlib import get_colors

# Map values to colors
values = np.array([1, 5, 10, 15, 20])
cmap = plt.cm.viridis
colors = get_colors(values, cmap, vmin=0, vmax=20, to255=False)
print("Colors shape:", colors.shape)

# Use with scatter plot
x = np.random.rand(100)
y = np.random.rand(100)
z = np.random.rand(100) * 100
colors = get_colors(z, plt.cm.jet, to255=False)
plt.scatter(x, y, c=colors)
plt.show()

# Handle NaN values
values_with_nan = np.array([1, 2, np.nan, 4, 5])
colors_nan = get_colors(values_with_nan, plt.cm.coolwarm, 
                         nancolor=[1, 1, 1, 1])  # White for NaN
\end{lstlisting}

\subsection{Function: \texttt{shiftedColorMap(cmap, start=0, midpoint=0.5, stop=1.0, name='shiftedcmap')}}

\subsubsection{Description}
Shift the center of a colormap to a specific data value. Useful for data with asymmetric ranges (e.g., $-15$ to $+5$) where you want the colormap center (typically white or neutral color) at zero.

\subsubsection{Input}
\begin{itemize}[leftmargin=*]
    \item \texttt{cmap}: matplotlib colormap -- The colormap to shift
    \item \texttt{start}: float -- Offset from lowest point [0.0, midpoint] (default: 0.0)
    \item \texttt{midpoint}: float -- New center position [0.0, 1.0] (default: 0.5)
    \begin{itemize}
        \item Calculate as: $1 - \frac{v_{\max}}{v_{\max} + |v_{\min}|}$
    \end{itemize}
    \item \texttt{stop}: float -- Offset from highest point [midpoint, 1.0] (default: 1.0)
    \item \texttt{name}: str -- Name for new colormap (default: \texttt{'shiftedcmap'})
\end{itemize}

\subsubsection{Output}
\begin{itemize}[leftmargin=*]
    \item \texttt{newcmap}: matplotlib.colors.LinearSegmentedColormap -- Shifted colormap
\end{itemize}

\subsubsection{Usage Example}
\begin{lstlisting}[language=Python, caption=Shifting colormap center for asymmetric data]
import matplotlib.pyplot as plt
import numpy as np
from plotlib import shiftedColorMap

# Data ranging from -15 to +5
data = np.random.randn(100, 100) * 10 - 5
vmin, vmax = -15, 5

# Calculate midpoint to center at zero
midpoint = 1 - vmax / (vmax + abs(vmin))
print(f"Midpoint: {midpoint:.3f}")  # 0.75

# Create shifted colormap
original_cmap = plt.cm.RdBu_r
shifted_cmap = shiftedColorMap(original_cmap, midpoint=midpoint)

# Plot with original vs shifted
fig, (ax1, ax2) = plt.subplots(1, 2, figsize=(12, 5))

im1 = ax1.imshow(data, cmap=original_cmap, vmin=vmin, vmax=vmax)
ax1.set_title('Original colormap')
plt.colorbar(im1, ax=ax1)

im2 = ax2.imshow(data, cmap=shifted_cmap, vmin=vmin, vmax=vmax)
ax2.set_title('Shifted colormap (zero = white)')
plt.colorbar(im2, ax=ax2)
plt.show()
\end{lstlisting}

\section{Best Practices}

\subsection{Figure Quality}
\begin{itemize}[leftmargin=*]
    \item Use DPI $\geq$ 300 for publication-quality figures
    \item Save in vector formats (PDF, SVG) when possible for scalability
    \item Enable cropping to remove unnecessary whitespace
    \item Use \texttt{tight\_layout=True} for multi-panel figures
\end{itemize}

\subsection{Color Scales}
\begin{itemize}[leftmargin=*]
    \item Use perceptually uniform colormaps (viridis, plasma) for continuous data
    \item Use diverging colormaps (RdBu) for data centered around zero
    \item Add inequality symbols ($\geq$, $\leq$) for saturated color scales
    \item Ensure sufficient color resolution (\texttt{cmapbins} $\geq$ 100)
\end{itemize}

\subsection{Projections}
\begin{itemize}[leftmargin=*]
    \item Use equal-area (Schmidt) for texture analysis and density plots
    \item Use equal-angle (Wulff) for angular relationships
    \item Use hemisphere projections for single-phase materials
    \item Use stereographic triangle for cubic crystal systems
\end{itemize}

\subsection{Performance}
\begin{itemize}[leftmargin=*]
    \item Pre-compute orientation data before plotting
    \item Use lower resolution (\texttt{nump}) for preview, higher for final output
    \item Cache symmetry operations for repeated calculations
    \item Batch export multiple formats in single call to \texttt{figsave()}
\end{itemize}

\section{Mathematical Background}

\subsection{Stereographic Projection}

The stereographic projection maps points from a sphere to a plane. For a point $\mathbf{p} = (x, y, z)$ on the unit sphere, the stereographic projection from the north pole $(0, 0, 1)$ to the $z = 0$ plane is:

\begin{equation}
(X, Y) = \left(\frac{x}{1-z}, \frac{y}{1-z}\right)
\end{equation}

For the equal-area (Schmidt) projection:
\begin{equation}
(X, Y) = \left(\frac{x}{\sqrt{1+z}}, \frac{y}{\sqrt{1+z}}\right)
\end{equation}

\subsection{Color Scale Normalization}

Values are normalized to colormap indices using:
\begin{equation}
v_n = \frac{v - v_{\min}}{v_{\max} - v_{\min}} \cdot (c_{\max} - c_{\min}) + c_{\min}
\end{equation}

where $v$ is the data value, $v_{\min}$ and $v_{\max}$ are the data range, and $c_{\min}$ and $c_{\max}$ are the colormap index range.

\subsection{Colormap Shifting}

For asymmetric data ranges, the midpoint $m$ should be calculated as:
\begin{equation}
m = 1 - \frac{v_{\max}}{v_{\max} + |v_{\min}|}
\end{equation}

For example, with data ranging from $-15$ to $+5$:
\begin{equation}
m = 1 - \frac{5}{5 + 15} = 1 - \frac{5}{20} = 0.75
\end{equation}

\section{Troubleshooting}

\subsection{Common Issues}

\subsubsection{Issue: Figures not displaying}
\begin{itemize}[leftmargin=*]
    \item \textbf{Solution}: Call \texttt{plt.show()} after creating the plot
    \item \textbf{Alternative}: Use \texttt{withdraw=False} parameter
\end{itemize}

\subsubsection{Issue: Low resolution output}
\begin{itemize}[leftmargin=*]
    \item \textbf{Solution}: Increase DPI: \texttt{p.figparam['dpi'] = 600}
    \item \textbf{Alternative}: Save as vector format (PDF, SVG)
\end{itemize}

\subsubsection{Issue: Colormap appears banded}
\begin{itemize}[leftmargin=*]
    \item \textbf{Solution}: Increase \texttt{cmapbins}: \texttt{getScales(cmapbins=1000)}
    \item \textbf{Alternative}: Use \texttt{cmapbinsmult} for automatic scaling
\end{itemize}

\subsubsection{Issue: Cropping not working}
\begin{itemize}[leftmargin=*]
    \item \textbf{Solution}: Install wand library: \texttt{pip install Wand}
    \item \textbf{Note}: Cropping only works for PNG format
\end{itemize}

\section{References}

\begin{itemize}[leftmargin=*]
    \item Stereographic projection: Kocks, U. F., Tomé, C. N., \& Wenk, H. R. (2000). \textit{Texture and anisotropy}. Cambridge University Press.
    \item Crystallographic conventions: Bunge, H. J. (2013). \textit{Texture analysis in materials science}. Elsevier.
    \item Color theory: Moreland, K. (2016). Why we use bad color maps and what you can do about it. \textit{Electronic Imaging}, 2016(16), 1-6.
\end{itemize}

\end{document}
